\subsubsection{Label Smoothing and Early Stopping Results}
\label{sec:lsearly}

% Owner: Member 5. Bullet-level skeleton; expand to ~1 page.

\paragraph{Label smoothing results.}
\draftnote{Expand this subsection into prose using the checklist below.}
\begin{itemize}
  \item Reference the method definition in \cref{sec:method:lsearly}
        and note that all shared training details follow
        \cref{sec:exp:training,sec:exp:eval}.
  \item Report best val acc / test acc / train--val gap / train time
        in a table.
  \item Briefly discuss confidence calibration (e.g. softmax max
        probability or ECE), since this is where label smoothing
        \citep{szegedy2016rethinking,muller2019labelsmoothing}
        differs most from other regularizers.
\end{itemize}

\paragraph{Early stopping results.}
\draftnote{Expand this subsection into prose using the checklist below.}
\begin{itemize}
  \item Reference the method definition in \cref{sec:method:lsearly}
        and note the shared protocol in
        \cref{sec:exp:training,sec:exp:eval}.
  \item Compare against fixed 30-epoch training: at which epoch does
        early stopping trigger? How much compute saved?
  \item Test accuracy at the early-stopped checkpoint vs final-epoch
        checkpoint.
\end{itemize}

\paragraph{Discussion.}
\label{sec:lsearly:discussion}
\draftnote{Answer the three Member 5 questions in short paragraphs.}
\begin{enumerate}
  \item Does label smoothing improve generalization here?
  \item Does early stopping actually help, or is it just ``stop
        early''?
  \item How do these output- and process-level regularizers relate to
        the parameter-level (\cref{sec:wd,sec:dropout}) and
        data-level (\cref{sec:augmentation}) ones?
\end{enumerate}
